%!TEX root = ../csuthesis_main.tex
% 设置中文摘要
\keywordscn{量子机器学习；音频分类；量子时序编码；量子时频分析神经网络；知识蒸馏}
\categorycn{TP183}
\itemcountcn{图 \totalfigures\ 幅，表 \totaltables\ 个，参考文献 \total{citnum}\ 篇}

\begin{abstractcn}\setlength{\baselineskip}{20pt}%\renewcommand{\baselinestretch}{1.0}
	音频分类是人工智能与数字信号处理领域的重要研究方向。近年来，经典深度学习模型虽取得显著进展，但也面临参数量庞大和算力瓶颈的挑战。随着量子计算的发展，量子音频处理应运而生，能够利用叠加与纠缠特性，以极低参数规模实现特征提取的指数级加速。然而，现有的量子音频处理方法往往忽略了音频信号的时序演化特性，且在含噪中等规模量子（NISQ）设备上面临深层线路训练困难等问题。因此，探索一种能高效表征时序特征并具备良好收敛性的量子神经网络架构是亟待解决的关键课题。本文提出了一套基于时序编码的量子神经网络音频分类与优化框架，从量子时频分析网络构建以及跨计算范式训练策略两个维度进行了系统性探索。本文的主要研究成果如下：
	
	（1）提出了一种基于量子时序编码的高效量子时频分析神经网络模型。针对传统量子编码难以捕捉音频动态时序特征的局限，该模型设计了基于双量子比特序列纠缠的量子时序编码方法，将离散化后的梅尔频率倒谱系数（MFCC）与音频时间索引进行全局纠缠映射，有效解决了时序特征向高维量子空间映射的难题。在此编码基础上，结合硬件高效拟设，构建了包含量子卷积层、量子池化层及量子测量模块的端到端网络架构。在 GTZAN 数据集的实验结果表明，该模型仅需 33 个固定的可训练参数便达到了 75\% 的分类准确率，在参数效率和特征捕获能力上展现了显著优势。
	
	（2）提出了“经典教师-量子学生”的跨计算范式知识蒸馏优化策略。为缓解变分量子线路在 NISQ 设备上易陷入局部最优及“贫瘠高原”现象的问题，本策略利用预训练的经典卷积神经网络模型作为教师网络，提取高维声学特征的连续概率分布。在此基础上，引入温度参数，构建了结合硬交叉熵与软 KL 散度的联合损失函数，反向引导参数量较少的量子线路进行参数更新。该软标签机制有助于平滑量子演化过程中的损失地形。实验结果表明，该蒸馏机制有效缩小了纯量子模型的泛化差距；模型在测试集上的准确率提升至 76.17\%，召回率达到 82.67\%，曲线下面积（AUC）为 0.8319，且 Brier 分数有所降低。这表明该策略在改善量子线路的收敛速度、预测可靠性及抗噪鲁棒性方面具有有效性。
	
\end{abstractcn}
